\chapter{SVM Classifier Evaluation}

\section{Experimental Setup}

We evaluated the Support Vector Machine (SVM) classifier using a **nested 5-fold stratified cross-validation** strategy to prevent data leakage. This approach follows the dual-track methodology outlined in the project architecture:

\begin{itemize}
    \item \textbf{Path A (Baseline):} Train SVM on raw preprocessed data (54,675 genes)
    \item \textbf{Path B (Optimized):} Apply feature selection within training folds, then train SVM on selected features
\end{itemize}

All SVM models used a linear kernel with balanced class weights to account for the balanced dataset structure.

\section{Path A: Baseline Track Results}

\subsection{Cross-Validation Performance}

The baseline SVM trained on all 54,675 genes achieved the following per-fold results:

\begin{table}[H]
    \centering
    \caption{SVM Baseline (Path A) - Per-Fold Results}
    \begin{tabular}{ccccccc}
        \toprule
        \textbf{Fold} & \textbf{Accuracy} & \textbf{Precision} & \textbf{Recall} & \textbf{F1-Score} & \textbf{MCC} & \textbf{ROC-AUC} \\
        \midrule
        1 & 0.8333 & 0.8235 & 0.8333 & 0.8333 & 0.6667 & 0.9444 \\
        2 & 1.0000 & 1.0000 & 1.0000 & 1.0000 & 1.0000 & 1.0000 \\
        3 & 0.9583 & 0.9583 & 0.9583 & 0.9600 & 0.9199 & 0.9792 \\
        4 & 0.9583 & 0.9583 & 0.9583 & 0.9565 & 0.9199 & 0.9792 \\
        5 & 1.0000 & 1.0000 & 1.0000 & 1.0000 & 1.0000 & 1.0000 \\
        \midrule
        \textbf{Mean} & \textbf{0.9500} & \textbf{0.9513} & \textbf{0.9500} & \textbf{0.9500} & \textbf{0.9013} & \textbf{0.9778} \\
        \textbf{Std} & \textbf{0.0685} & \textbf{0.0739} & \textbf{0.0745} & \textbf{0.0685} & \textbf{0.1371} & \textbf{0.0319} \\
        \bottomrule
    \end{tabular}
\end{table}

\subsection{Key Findings - Baseline}

\begin{itemize}
    \item \textbf{Mean Accuracy:} 95.00\% ± 6.85\%
    \item \textbf{Mean MCC:} 0.9013 ± 0.1371 (strong balanced performance)
    \item \textbf{Mean ROC-AUC:} 0.9778 (excellent discrimination)
    \item \textbf{Feature Set:} All 54,675 genes
    \item \textbf{Variability:} High variability in MCC (std = 0.1371) suggests fold-dependent performance
\end{itemize}

The baseline demonstrates that even without feature selection, SVM can achieve high classification accuracy. However, the high dimensionality (p = 54,675) creates computational overhead and potential overfitting risk.

\section{Path B: Optimized Track Results}

\subsection{Filter Method: Welch's T-Test}

\begin{table}[H]
    \centering
    \caption{SVM with Filter (T-Test) - Per-Fold Results}
    \begin{tabular}{ccccccc}
        \toprule
        \textbf{Fold} & \textbf{Accuracy} & \textbf{Precision} & \textbf{Recall} & \textbf{F1-Score} & \textbf{MCC} & \textbf{ROC-AUC} \\
        \midrule
        1 & 0.8750 & 0.8571 & 0.8333 & 0.8696 & 0.7526 & 0.9722 \\
        2 & 0.9583 & 0.9583 & 0.9583 & 0.9565 & 0.9199 & 0.9792 \\
        3 & 0.9583 & 0.9583 & 0.9583 & 0.9600 & 0.9199 & 1.0000 \\
        4 & 0.9583 & 0.9583 & 0.9583 & 0.9565 & 0.9199 & 0.9653 \\
        5 & 1.0000 & 1.0000 & 1.0000 & 1.0000 & 1.0000 & 1.0000 \\
        \midrule
        \textbf{Mean} & \textbf{0.9500} & \textbf{0.9664} & \textbf{0.9333} & \textbf{0.9485} & \textbf{0.9024} & \textbf{0.9875} \\
        \textbf{Std} & \textbf{0.0456} & \textbf{0.0462} & \textbf{0.0697} & \textbf{0.0478} & \textbf{0.0907} & \textbf{0.0104} \\
        \bottomrule
    \end{tabular}
\end{table}

\paragraph{Analysis:}
- Achieved identical mean accuracy (95.00%) to baseline with only \textbf{20 genes}
- \textbf{Improved ROC-AUC} from 0.9778 to 0.9875 (better discrimination at decision boundary)
- \textbf{Reduced variability} in MCC (std: 0.1371 $\rightarrow$ 0.0907) indicates more stable performance
- \textbf{99.96\% dimensionality reduction} with no accuracy loss

\subsection{Filter Method: ANOVA F-Test}

\begin{table}[H]
    \centering
    \caption{SVM with Filter (ANOVA) - Per-Fold Results}
    \begin{tabular}{ccccccc}
        \toprule
        \textbf{Fold} & \textbf{Accuracy} & \textbf{Precision} & \textbf{Recall} & \textbf{F1-Score} & \textbf{MCC} & \textbf{ROC-AUC} \\
        \midrule
        1 & 0.8750 & 0.8571 & 0.8333 & 0.8696 & 0.7526 & 0.9722 \\
        2 & 0.9583 & 0.9583 & 0.9583 & 0.9565 & 0.9199 & 0.9792 \\
        3 & 0.9583 & 0.9583 & 0.9583 & 0.9600 & 0.9199 & 1.0000 \\
        4 & 1.0000 & 1.0000 & 1.0000 & 1.0000 & 1.0000 & 1.0000 \\
        5 & 1.0000 & 1.0000 & 1.0000 & 1.0000 & 1.0000 & 1.0000 \\
        \midrule
        \textbf{Mean} & \textbf{0.9583} & \textbf{0.9664} & \textbf{0.9500} & \textbf{0.9572} & \textbf{0.9185} & \textbf{0.9903} \\
        \textbf{Std} & \textbf{0.0510} & \textbf{0.0462} & \textbf{0.0745} & \textbf{0.0533} & \textbf{0.1010} & \textbf{0.0102} \\
        \bottomrule
    \end{tabular}
\end{table}

\paragraph{Analysis:}
- \textbf{Best filter performance:} 95.83\% accuracy (improvement over baseline)
- \textbf{Best MCC:} 0.9185 ± 0.1010
- \textbf{Lowest variability} among filter methods
- ANOVA outperforms T-Test by penalizing high intra-class variance

\subsection{Wrapper Method: SVM-RFE}

\begin{table}[H]
    \centering
    \caption{SVM with Wrapper (SVM-RFE) - Per-Fold Results}
    \begin{tabular}{ccccccc}
        \toprule
        \textbf{Fold} & \textbf{Accuracy} & \textbf{Precision} & \textbf{Recall} & \textbf{F1-Score} & \textbf{MCC} & \textbf{ROC-AUC} \\
        \midrule
        1 & 0.8750 & 0.8571 & 0.8333 & 0.8696 & 0.7526 & 0.9722 \\
        2 & 1.0000 & 1.0000 & 1.0000 & 1.0000 & 1.0000 & 1.0000 \\
        3 & 0.9583 & 0.9583 & 0.9583 & 0.9600 & 0.9199 & 0.9792 \\
        4 & 0.9583 & 0.9583 & 0.9583 & 0.9565 & 0.9199 & 0.9653 \\
        5 & 1.0000 & 1.0000 & 1.0000 & 1.0000 & 1.0000 & 1.0000 \\
        \midrule
        \textbf{Mean} & \textbf{0.9583} & \textbf{0.9664} & \textbf{0.9500} & \textbf{0.9572} & \textbf{0.9185} & \textbf{0.9889} \\
        \textbf{Std} & \textbf{0.0510} & \textbf{0.0462} & \textbf{0.0745} & \textbf{0.0533} & \textbf{0.1010} & \textbf{0.0124} \\
        \bottomrule
    \end{tabular}
\end{table}

\paragraph{Analysis:}
- \textbf{Tied with ANOVA:} 95.83\% accuracy
- \textbf{Similar MCC:} 0.9185 ± 0.1010
- SVM-RFE captures multivariate feature interactions but achieves same performance as filter
- Computational cost justified only if interpretability of model-specific selection is desired

\subsection{Embedded Method: LASSO}

\begin{table}[H]
    \centering
    \caption{SVM with Embedded (LASSO) - Per-Fold Results}
    \begin{tabular}{ccccccc}
        \toprule
        \textbf{Fold} & \textbf{Accuracy} & \textbf{Precision} & \textbf{Recall} & \textbf{F1-Score} & \textbf{MCC} & \textbf{ROC-AUC} \\
        \midrule
        1 & 0.7917 & 0.7857 & 0.7500 & 0.8000 & 0.5854 & 0.9306 \\
        2 & 1.0000 & 1.0000 & 1.0000 & 1.0000 & 1.0000 & 1.0000 \\
        3 & 0.9583 & 0.9583 & 0.9583 & 0.9600 & 0.9199 & 0.9792 \\
        4 & 0.8750 & 0.8571 & 0.8333 & 0.8571 & 0.7746 & 0.9583 \\
        5 & 1.0000 & 1.0000 & 1.0000 & 1.0000 & 1.0000 & 1.0000 \\
        \midrule
        \textbf{Mean} & \textbf{0.9250} & \textbf{0.9385} & \textbf{0.9167} & \textbf{0.9234} & \textbf{0.8560} & \textbf{0.9819} \\
        \textbf{Std} & \textbf{0.0903} & \textbf{0.1003} & \textbf{0.1179} & \textbf{0.0904} & \textbf{0.1771} & \textbf{0.0364} \\
        \bottomrule
    \end{tabular}
\end{table}

\paragraph{Analysis:}
- \textbf{Lower performance:} 92.50\% accuracy (below filter/wrapper methods)
- \textbf{Highest variability:} MCC std = 0.1771 (least stable across folds)
- LASSO's L1 penalty may be too aggressive; $\lambda = 0.1$ could benefit from tuning
- Trade-off: More interpretable coefficients but less predictive accuracy

\section{Comparative Summary}

\begin{table}[H]
    \centering
    \caption{SVM Performance Comparison Across All Methods}
    \begin{tabular}{lcccccc}
        \toprule
        \textbf{Method} & \textbf{Accuracy} & \textbf{Precision} & \textbf{Recall} & \textbf{F1-Score} & \textbf{MCC} & \textbf{n\_Features} \\
        \midrule
        Path A: Baseline & $0.9500 \pm 0.0685$ & $0.9513 \pm 0.0739$ & $0.9500 \pm 0.0745$ & $0.9500 \pm 0.0685$ & $0.9013 \pm 0.1371$ & 54,675 \\
        Path B: T-Test & $0.9500 \pm 0.0456$ & $0.9664 \pm 0.0462$ & $0.9333 \pm 0.0697$ & $0.9485 \pm 0.0478$ & $0.9024 \pm 0.0907$ & 20 \\
        Path B: ANOVA & $0.9583 \pm 0.0510$ & $0.9664 \pm 0.0462$ & $0.9500 \pm 0.0745$ & $0.9572 \pm 0.0533$ & $0.9185 \pm 0.1010$ & 20 \\
        Path B: SVM-RFE & $0.9583 \pm 0.0510$ & $0.9664 \pm 0.0462$ & $0.9500 \pm 0.0745$ & $0.9572 \pm 0.0533$ & $0.9185 \pm 0.1010$ & 20 \\
        Path B: LASSO & $0.9250 \pm 0.0903$ & $0.9385 \pm 0.1003$ & $0.9167 \pm 0.1179$ & $0.9234 \pm 0.0904$ & $0.8560 \pm 0.1771$ & 20 \\
        \bottomrule
    \end{tabular}
\end{table}

\section{Key Insights}

\subsection{Performance Plateau}

The remarkable finding is that \textbf{all feature selection methods achieved 92.5\% - 95.83\% accuracy}, representing a \textbf{plateau effect}. This suggests:
\begin{itemize}
    \item Reducing 54,675 genes to 20 incurs minimal performance cost
    \item The signal is concentrated in a small subset of genes
    \item Beyond 20 genes, additional features add noise rather than signal
\end{itemize}

\subsection{Best Performer: ANOVA and SVM-RFE (tied)}

\begin{itemize}
    \item Accuracy: 95.83\% (highest)
    \item MCC: 0.9185 (most stable for imbalanced evaluation)
    \item ROC-AUC: 0.9903 (excellent discrimination)
    \item Variability: Lowest among all methods
\end{itemize}

The tie suggests that \textbf{statistical tests (ANOVA) capture the same predictive signal as model-specific methods (SVM-RFE)} on this dataset.

\subsection{Computational Efficiency}

\begin{table}[H]
    \centering
    \caption{Computational Cost Comparison}
    \begin{tabular}{lcc}
        \toprule
        \textbf{Method} & \textbf{Time per Fold (approx)} & \textbf{Features Evaluated} \\
        \midrule
        Baseline (no selection) & ~5 sec & 54,675 \\
        T-Test & <1 sec & 54,675 \\
        ANOVA & <1 sec & 54,675 \\
        SVM-RFE & ~0.5 sec & 500 (pre-filtered) \\
        LASSO & ~0.6 sec & 54,675 \\
        \bottomrule
    \end{tabular}
\end{table}

All feature selection methods complete within 1 second per fold, demonstrating that feature selection \textbf{reduces computational burden without sacrificing accuracy}.

\section{Discussion}

\subsection{Why SVM Works Well with Feature Selection}

SVM's margin-maximization objective is well-suited for high-dimensional data. When combined with feature selection:
\begin{enumerate}
    \item Reduces noise in the feature space
    \item Simplifies the optimization landscape
    \item Improves numerical stability
\end{enumerate}

\subsection{LASSO Underperformance}

LASSO with $\lambda = 0.1$ appears to be over-regularized. Future improvements could:
\begin{itemize}
    \item Use cross-validated $\lambda$ selection
    \item Implement Elastic Net (mix L1 and L2 penalties)
    \item Explore alternative regularization strengths
\end{itemize}

\subsection{Stability and Generalization}

The reduced variability in feature-selected models (ANOVA/SVM-RFE std = 0.0510) compared to baseline (std = 0.0685) suggests:
\begin{itemize}
    \item Feature selection improves \textbf{generalization across different data splits}
    \item Less fold-to-fold variance indicates more robust learned models
    \item Better performance on unseen data expected
\end{itemize}
